\documentclass{article}
\usepackage{tikz}
\usepackage{amsmath}
\usepackage{float}
\usepackage{xcolor}
\usepackage{cite}
\usepackage{graphicx} 
\usepackage{amsmath}
\usepackage{amssymb}
\usepackage{amsfonts}
\usepackage[a4paper, margin=2cm]{geometry}

\title{Sampling of the Keys for a PRF}
\author{Jocelyn Fagard}
\date{March 2026}

\begin{document}

\maketitle

\section{Context}
Let $n,m \in \mathbb{N}^*$; let $q$ be a power of a prime integer and consider the ring
\[
A := \mathbb{F}_q[X]/(X^m + 1)
\]
as well as $(n+1)$ polynomials $K:=(k_0, \dots, k_n) \in A^{n+1}$.
We also assume we have a rounding function $S : \mathbb{F}_q^m \to \{0,1\}^m$ which is the concatenation of a function $s : \mathbb{F}_q \to \{0,1\}$.

We are then interested in studying the function
\[
F_K : \{0,1\}^n \to \{0,1\}^m
\]
that maps $x := (x_1, \dots, x_n)$ to
\[
S\!\left( k_0 \prod_{i=1}^n k_i^{x_i} \right).
\] that we want to be pseudorandom.

\section{Shifted keys}
Let $k_0 \in A$, here we assume that for all $i \in \{0,\dots,n-1\}$ $k_{i+1}:= \sigma(k_i)$ where $\sigma$ is the cyclic shift  : $P \in \mathbb{F}_q[X]_{<m} \to X  P$ mod $(X^m-1)$.
Given a polynomial $r$ and an non-negative integer $l$ we note $r_l$ the $(l+1)$-th coefficient of $r$. Let's start we a remark.

\textbf{Claim.} Assume that $m$ is even.
Let $a,b \in A$. Then we have the following equality:
\[
(ab)_{\frac{m}{2} - 1}
=
- \left( \sigma^{\frac{m}{2}}(a)\, \sigma^{-\frac{m}{2}}(b) \right)_{\frac{m}{2} - 1}.
\]

\medskip

\textbf{Proof.}
This can be obtained by direct computation. We write
\[
(ab)_{\frac{m}{2} - 1}
=
\sum_{i=0}^{\frac{m}{2}-1} a_i \, b_{\frac{m}{2}-1 - i}
-
\sum_{i=\frac{m}{2}}^{m-1} a_i \, b_{\frac{m}{2}-1 - i},
\]
and
\[
\left( \sigma^{\frac{m}{2}}(a)\, \sigma^{-\frac{m}{2}}(b) \right)_{\frac{m}{2} - 1}
=
\sum_{i=0}^{\frac{m}{2}-1} a_{i-\frac{m}{2}} \, b_{\frac{m}{2}-1-i+\frac{m}{2}}
-
\sum_{i=\frac{m}{2}}^{m-1} a_{i-\frac{m}{2}} \, b_{\frac{m}{2}-1-i+\frac{m}{2}},
\]
where indices are taken modulo $m$.

By a change of variable, we obtain
\[
\left( \sigma^{\frac{m}{2}}(a)\, \sigma^{-\frac{m}{2}}(b) \right)_{\frac{m}{2} - 1}
=
\sum_{i=\frac{m}{2}}^{m-1} a_i \, b_{m-1-i}
-
\sum_{i=0}^{\frac{m}{2}-1} a_i \, b_{\frac{m}{2}-1-i}
=
- (ab)_{\frac{m}{2} - 1}.
\]
\hfill $\square$

\textbf{Consequence.}
Assume that $m$ is even and that $n \geq \frac{m}{2}$. Then, by the previous claim, we have
\[
(k_0 k_{\frac{m}{2}})_{\frac{m}{2}-1}
=
- \left( \sigma^{\frac{m}{2}}(k_0)\, \sigma^{-\frac{m}{2}}(k_{\frac{m}{2}}) \right)_{\frac{m}{2}-1}.
\]
Moreover, by the way the keys are generated, we also have
\[
 \left( \sigma^{\frac{m}{2}}(k_0)\, \sigma^{-\frac{m}{2}}(k_{\frac{m}{2}}) \right)_{\frac{m}{2}-1}
=
 (k_{\frac{m}{2}} k_0)_{\frac{m}{2}-1}.
\]
Hence,
\[
(k_0 k_{\frac{m}{2}})_{\frac{m}{2}-1} = 0.
\]

\medskip

Assume now that $s(0) = b \in \{0,1\}$. Then a simple algorithm to distinguish $F_K$ from a random function is the following:

\begin{itemize}
    \item Query the function on $x := (\delta_{i, \frac{m}{2}})_{1 \leq i \leq n}$.
    \item Let $y := (y_1, \dots, y_m)$ be the output.
    \item If $y_{\frac{m}{2}-1} = b$, output that the function is $F_K$; otherwise, output that it is random function.
\end{itemize}

The success probability of this distinguisher is
\[
\frac{1}{2} + \frac{1}{4} = \frac{3}{4}.
\]

\section{Linear relation between keys}
In this section, we assume that the keys $k_1, k_2, \dots, k_n$ are sampled independently and uniformly at random from a $d$-dimensional linear subspace of $\mathbb{F}_q[X]_{<m}$, with we identify with $\mathbb{F}_q^m$, where $1 \leq d \leq m$. We denote by $(e_1, \dots, e_d)$ a basis of this subspace.

We also fix an integer $1 \leq \ell \leq n$, which represents the Hamming weight $w_H(x)$ of the input $x := (x_1, x_2, \dots, x_n)$ used in our queries.

\medskip

Let
\[
E_\ell := \mathrm{span} \left\{ k_0 \prod_{i=1}^n k_i^{x_i} \;\middle|\; x \in \{0,1\}^n,\; w_H(x) = \ell \right\}.
\]
Then we have the inclusion
\[
E_\ell \subseteq \mathrm{span} \left\{ \prod_{i=1}^d e_i^{z_i} \;\middle|\; z_i \in \mathbb{N},\; \sum_{i=1}^d z_i = \ell \right\}.
\] Therefore, $E_\ell$ is generated by at most $\binom{d+\ell-1}{d -1}$ (which is the number of monomials of degree $\ell$ in d variables) elements. This provides an upper bound on its dimension, denoted $c_{d,\ell}$.

Now assume we take $x, y \in \{0,1\}^n$ uniformly at random among messages of Hamming weight $\ell$ with $x \neq y$.

We can then expect
\[
\Pr[F_K(x) = F_K(y)] = \frac{1}{q^{c_{d,\ell}}} \geq \frac{1}{q^{\binom{d+\ell-1}{d-1}}} = \frac{1}{2^{\log_2(q)\binom{d+\ell-1}{d-1}}} =: p_{d,\ell}.
\]
This probability is maximized for $\ell=1$, where it equals
\[
p_{d,1} = \frac{1}{2^{\log_2(q)\, d}}.
\]
Note that for a truly random function, this collision probability is $2^{-m}$.

Let's now present a Distinguisher.
\medskip

\textbf{Distinguisher:}
\begin{itemize}
    \item Take two different vectors of Hamming weight 1 and query the function on them.
    \item If there is a collision, output that the function is $F_K$; otherwise, output that it is either $F_K$ or a random function with probability $1/2$ each.
\end{itemize}

\medskip

The success probability of this distinguisher can then be expressed as
\[
\Pr[\text{success}] \geq \frac{1}{2} \cdot \frac{1}{2} \cdot \left(1 - 2^{-m}\right) 
+ \frac{1}{2} \left( p_{d,1} + (1 - p_{d,1}) \cdot \frac{1}{2} \right) = \frac{1}{2} + \frac{p_{d,1}}{4} - \frac{1}{2^{m+2}}.
\]
\textbf{Consequence.}  
The function $F_K$ is not pseudorandom if $p_{d,1}$ is the inverse of a polynomial in $m$, that is, if 
\[
d = \log_2 \bigl( U(m) \bigr)
\]
for some polynomial $U \in \mathbb{R}[X]$.

\textbf{Remark.}  
Fortunately, the probability that the keys lie in a low-dimensional space when they are sampled uniformly at random and independently is very low, since
\[
\Pr\bigl[\dim(\mathrm{span}(k_1, k_2, \dots, k_n)) = d \bigr]
= \frac{\prod_{i=0}^{d-1} \frac{(q^m - q^i)(q^n - q^i)}{q^d - q^i}}{q^{mn}}.
\]

\section{Freshman Dream}

In this section, we remark that no nontrivial ring $R$ exists such that for all $a,b,c,d \in R$, the equality
\[
(a+b)(c+d) = ac + bd
\]
holds.

To see this, suppose such a ring exists. Let $a \in R$, and denote by $0$ and $1$ the additive and multiplicative identities of $R$, respectively.

On the one hand, using the assumed property, we have
\[
(a+0)(1+0) = a \cdot 1 = a.
\]

On the other hand, using the assumed property and commutativity of addition, we get
\[
(a+0)(1+0) = (0+a)(1+0) = 0 \cdot 1 + a \cdot 0 = 0 + 0 = 0.
\]

Hence, $a = 0$ for all $a \in R$, which means that $R$ is the zero ring. Therefore, no nontrivial ring satisfies the given property.

\section{Some open problems about polynomial permutations}

Let $p$ be a prime number and let $d \geq 2$.

A polynomial permutation of $\mathbb{Z}_{p^d}$ is a polynomial 
$f \in \mathbb{Z}_{p^d}[X]$ that induces a permutation of the set $\mathbb{Z}_{p^d}$.

In previous works, it was proved that:

\vspace{0.3cm}

\textbf{Theorem.} 

A polynomial $f \in \mathbb{Z}_{p^d}[X]$ is a permutation polynomial 
if and only if its reduction modulo $p$ is a permutation polynomial of $\mathbb{Z}_p$ 
and $f'$ does not vanish on $\mathbb{Z}_p$.

\vspace{0.3cm}

In addition, the number of permutation polynomials of degree $n$ over $\mathbb{Z}_{p^d}$, 
denoted by $N(p^d,n)$, was left as an open question 
for $3 \leq n \leq 2p-2$.

Here, we address these questions.

\vspace{0.3cm}

\textbf{Lemma.} 

The number of permutation polynomials of $\mathbb{Z}_p$ 
of degree $p-1$ is zero.

Hence, for all $d \geq 2$, 
\[
N(p^d,p-1) = 0.
\]

\vspace{0.3cm}

\textbf{Proof.}

Given a permutation $\sigma$ of $\{0,1,\dots,p-1\}$, 
Lagrange interpolation shows that the unique polynomial 
of degree at most $p-1$ inducing this permutation is

\[
\sum_{i=0}^{p-1} \sigma(i) 
\prod_{\substack{0 \leq j \leq p-1 \\ j \neq i}} 
\frac{X - j}{i - j}.
\]

The coefficient of $X^{p-1}$ is therefore

\[
\sum_{i=0}^{p-1} \sigma(i) 
= \sum_{i=0}^{p-1} i 
= 0 \quad \text{in } \mathbb{Z}_p.
\]

Thus, no permutation polynomial over $\mathbb{Z}_p$ 
has degree exactly $p-1$.

\hfill $\square$

\vspace{0.3cm}
\textbf{Lemma.}

Let $p \geq 3$, we have:

\[
\frac{1}{2} \leq\frac{N(p, p-2)}{p!} \leq 1.
\]

That is at least half of polynomial permutations have degree $p-2$.

\textbf{Proof.}
Again, given a permutation $\sigma$ of $\mathbb{Z}_p$ using Lagrange interpolation formula, we have that the coefficient in front of the monom of degree $p-2$ is :

\[
- \sum_{i=0}^{p-1} \sigma(i) (\sum_{j \neq i} (-j)) = \sum_{i=0}^{p-1} \sigma(i)\cdot i = <id, \sigma>
\] 

where $<\cdot,\cdot>$ denotes the scalar product, $id = (0,1,\cdots,p-1)$ and $\sigma = (\sigma(0),\sigma(1),\cdots, \sigma(p-1))$. Hence $\sigma$ has degree $p-2$ if and only if it is not orthogonal to $id$.
\vspace{0.3cm}

Given $j \in \mathbb{Z}_p$ we define, $A_j :=\{\tau \text{ permutation} \mid <\tau, id> = j \}$ and $a_j:=\#A_j$. We first claim that $\forall i \neq j \in \mathbb{Z}_p^*$ , $a_i=a_j$ by remarking that $\forall \tau$ such that $<\tau,id> \neq 0$ there is exactly one permutation among the non zero co-linear permutations $\{ \lambda \cdot\tau \mid \lambda \in \mathbb{Z}_p^*\}$ in every $A_j$ for $j\neq0$.  
We denote this common value of these $a_i$'s with $i\neq 0$ by $b$.

\vspace{0.3cm}

Secondly, let's note $f_{i,j}: \tau \to (\tau_i,\tau_j) \tau$ the application that switches the i-th and the j-th values. Then we remark $<f_{i,j}(\tau), id> = <\tau,id> + (\tau(j) - \tau(i))(i-j) \neq <\tau,id>$ for $\tau$ a permutation. So, each element of some $A_l$ is associated to an element which is not in $A_l$. 
\vspace{0.3cm}

Now, we consider the graph where the vertices are the permutations and there is an edge between two permutations if one is obtained by applying $f_{0,1}$. In this graph every vertices are in couple because $f_{0,1} \circ f_{0,1} = id$. Each of the $a_0$ vertex in $A_0$ is associated to a distinct elements among the $(p-1)b$ which are not in $A_0$. So $(p-1)b \geq a_0$. And consequently the number of permutation polynomials of degree $p-2$ which is $(p-1)b$ represents at least half of permutations.

\vspace{0.3cm}

\hfill $\square$

\textbf{Proposition.}

Let $1 \leq n \leq p-1$ and $d \geq 2$. Then
\[
\frac{N(p^d,\leq p-1+n)}{p^{(d-1)(p+n)}}
= p! \left( 
p^n - \sum_{k=1}^{p} (-1)^{k-1} \binom{p}{k} \, p^{\max(n-k,0)}
\right).
\]

\vspace{0.3cm}

\textbf{Proof.}

Any polynomial $f \in \mathbb{Z}_{p}[X]$ of degree at most $p-1+n$ can be uniquely written in the form
\[
f(x) = (x^p - x) q(x) + \sigma(x),
\]
where $\sigma$ is a polynomial of degree at most $p-1$ and $q$ is a polynomial of degree at most $n-1$.

Moreover, $\sigma$ induces a permutation of $\mathbb{Z}_p$, and every such permutation arises uniquely in this way.

\vspace{0.2cm}

By the previous theorem, $f$ is a permutation polynomial over $\mathbb{Z}_{p^d}$ if and only if its derivative does not vanish on $\mathbb{Z}_p$. A direct computation shows that
\[
f'(x) \equiv -q(x) + \sigma'(x) ,
\]
so the condition becomes
\[
\forall i \in \mathbb{Z}_p, \quad \sigma'(i) \neq q(i).
\]

\vspace{0.2cm}

Fix a permutation polynomial $\sigma$. We want to count the number of polynomials $q \in \mathbb{Z}_p[X]$ of degree at most $n-1$ such that
\[
\forall i \in \mathbb{Z}_p, \quad q(i) \neq \sigma'(i).
\]

For each $i \in \mathbb{Z}_p$, define
\[
A_i := \left\{ q \in \mathbb{Z}_p[X]_{< n} \;\middle|\; q(i) = \sigma'(i) \right\}.
\]

We are therefore interested in the cardinality of the complement of $\bigcup_{i \in \mathbb{Z}_p} A_i$.

\vspace{0.2cm}

Using the inclusion--exclusion principle, and the fact that for any distinct 
$i_1,\dots,i_k \in \mathbb{Z}_p$, the intersection
\[
A_{i_1} \cap \cdots \cap A_{i_k}
\]
has cardinality $p^{n-k}$ if $1 \leq k \leq n$, and $1$ if $k \geq n$, we obtain
\[
\left|\bigcup_{i \in \mathbb{Z}_p} A_i\right|
= \sum_{k=1}^{p} (-1)^{k-1} \binom{p}{k} \, p^{\max(n-k,0)}.
\]

\vspace{0.2cm}

Hence, the number of admissible polynomials $q$ is
\[
p^n - \sum_{k=1}^{p} (-1)^{k-1} \binom{p}{k} \, p^{\max(n-k,0)}.
\]

\vspace{0.2cm}

Since there are $p!$ possible permutation polynomials $\sigma$ over $\mathbb{Z}_p$, we finally obtain
\[
N(p^d,\leq p-1+n)
= p! \left( 
p^n - \sum_{k=1}^{p} (-1)^{k-1} \binom{p}{k} \, p^{\max(n-k,0)}
\right).
\]

\hfill $\square$
\end{document}
